\textbf{1. Kiến thức cần nhớ}
\begin{itemize}
\item \textbf{Định nghĩa}: Đường tròn tâm $O$ bán kính $R$ ($R > 0$) là hình gồm các điểm cách $O$ một khoảng bằng $R$, kí hiệu là $(O; R)$.
\item \textbf{Vị trí tương đối của một điểm $M$ đối với $(O; R)$}:
\begin{itemize}
\item $M$ nằm trong $(O) \iff OM < R$.
\item $M$ nằm trên $(O) \iff OM = R$.
\item $M$ nằm ngoài $(O) \iff OM > R$.
\end{itemize}
\item \textbf{Tính đối xứng}:
\begin{itemize}
\item Đường tròn là hình có \textbf{tâm đối xứng}: Tâm $O$ là tâm đối xứng của đường tròn.
\item Đường tròn là hình có \textbf{trục đối xứng}: Bất kì đường kính nào cũng là một trục đối xứng của đường tròn.
\end{itemize}
\item \textbf{Sự xác định đường tròn}: Qua $3$ điểm không thẳng hàng, ta vẽ được \textbf{một và chỉ một} đường tròn. Tâm của đường tròn này là giao điểm của ba đường trung trực của tam giác tạo bởi ba điểm đó.
\item Tam giác vuông nội tiếp đường tròn có cạnh huyền là đường kính (tâm là trung điểm cạnh huyền).
\end{itemize}

\textbf{2. Phương pháp giải các dạng}

%% DANG 9H5?13-1
\textbf{Dạng 1. Vị trí của điểm đối với đường tròn}
\begin{enumerate}
\item Tính khoảng cách từ điểm $M$ đến tâm $O$ (sử dụng định lí Pythagore hoặc hệ thức tọa độ).
\item So sánh $OM$ với bán kính $R$:
\begin{itemize}
\item $OM < R \implies M$ nằm trong đường tròn.
\item $OM = R \implies M$ nằm trên đường tròn.
\item $OM > R \implies M$ nằm ngoài đường tròn.
\end{itemize}
\end{enumerate}

%% DANG 9H5?13-2
\textbf{Dạng 2. Tâm đối xứng, trục đối xứng của đường tròn}
\begin{enumerate}
\item Lấy đối xứng một điểm qua tâm $O$: nếu $A \in (O)$ thì điểm đối xứng $A'$ qua $O$ cũng thuộc $(O)$ và đoạn $AA'$ là đường kính.
\item Lấy đối xứng qua một đường kính: dây cung vuông góc với đường kính thì nhận đường kính làm trục đối xứng.
\end{enumerate}

%% DANG 9H5?13-3
\textbf{Dạng 3. Chứng minh các điểm cùng thuộc một đường tròn}
\begin{enumerate}
\item \textbf{Cách 1}: Chứng minh các điểm cùng cách đều một điểm cố định $O$ một khoảng không đổi $R$.
\item \textbf{Cách 2}: Chứng minh các tam giác vuông có chung cạnh huyền $\rightarrow$ các đỉnh góc vuông và hai đầu cạnh huyền cùng thuộc đường tròn đường kính là cạnh huyền chung đó.
\end{enumerate}
